%-----------------------------------------------------------------------
\section{How computer graphics builds caustics}
\label{sec:cg}
%-----------------------------------------------------------------------

Every caustic method in production use estimates \eqref{eq:irradiance}.
Rendering caustics by following light through curved bodies is as old as
global-illumination rendering itself \cite{mitchell1992}; what has changed since
is the estimator, not the object being estimated. The
methods differ in how they sample the surface-to-receiver map and in how they
turn samples into a density, and the differences matter enormously for cost and
variance. They do not differ in the integrand, and the integrand is what
Lemma~\ref{lem:diverge} says is unbounded. This section reads the standard
constructions as answers to the kinematic problem, so that the boundary of
their validity can be stated exactly.

\subsection{Estimating a push-forward}
\label{sec:two-families}

Estimating $I(\mathbf{u})$ requires samples of $\Hm$ and a way to convert them
into a density at $\mathbf{u}$. There are exactly two families in practice.

\emph{Forward} methods sample the surface, push the samples through $\Hm$, and
estimate the resulting density where the samples land. Photon mapping
\cite{jensen1996} and its progressive variant \cite{hachisuka2008} are the
canonical members. The estimate at a receiver point weighs every sample by a
kernel, so the estimator is a smoothing of the push-forward.

\emph{Backward} methods sample from the camera, follow the transport, and
connect to light-carrying vertices, weighting the connection by a geometric
term that contains the same $1/|\detJ|$ factor. Path tracing with next-event
estimation, bidirectional path tracing, and the manifold-based methods
\cite{hanika2015,zeltner2020} are members. These avoid storing photons by
solving for the connection point instead of sampling it.

Both families are Monte Carlo estimators of the same integral. Both are
unbiased for the \emph{integral} of $I$ over a pixel, which is what a renderer
must produce. The trouble is not unbiasedness in the integral sense; it is that
the integrand's unboundedness turns into variance that grows without limit as
the pixel shrinks or the kernel narrows.

\subsection{Forward: photon mapping and the bias at the fold}
\label{sec:photonmap}

Photon mapping deposits $N$ samples $\mathbf{u}_i = \Hm(\mathbf{x}_i)$ and
estimates the density with a kernel $K_h$ of bandwidth $h$,
\begin{equation}
\hat I(\mathbf{u}) \;=\; \frac{1}{N} \sum_{i=1}^{N}
K_h\!\big(\mathbf{u} - \mathbf{u}_i\big).
\label{eq:kde}
\end{equation}
When the $\mathbf{x}_i$ are uniform on the surface, $\hat I$ estimates the
push-forward density of \eqref{eq:irradiance} after smoothing. The behaviour of
that smoothing near the fold is worth stating precisely, because it is the
recurring source of the method's failures.

\begin{lemma}[Kernel smoothing cannot resolve the fold]
\label{lem:kdebias}
Let $K_h$ be a kernel with $\int K_h = 1$, finite second moment
$\sigma^2 = \int \mathbf{z}^2 K(\mathbf{z})\dd\mathbf{z}$, and let $I$ be the
push-forward density of \eqref{eq:irradiance}. Then
\begin{equation}
\mathbb{E}\big[\hat I(\mathbf{u})\big] - I(\mathbf{u})
\;=\; \frac{h^2}{2}\sigma^2 \nabla^2 I(\mathbf{u}) + \order{h^4},
\label{eq:kdebias}
\end{equation}
and near the fold, where $I \sim (2\sqrt{b})^{-1}$ in the normal-form coordinate
$b$ of \eqref{eq:normalform}, the second derivative behaves as
$\nabla^2 I \sim b^{-5/2}$ and the bias in \eqref{eq:kdebias} is unbounded for
every fixed $h > 0$.
\end{lemma}

\begin{proof}
Equation~\eqref{eq:kdebias} is the standard kernel-density bias expansion
\cite{silverman1986}, obtained by Taylor-expanding $I$ under the convolution.
For the fold, differentiating \eqref{eq:folddensity} twice in $b$ gives
$\rho''(b) = \tfrac{3}{8}b^{-5/2}$, which diverges as $b\to 0^+$. Since
\eqref{eq:kdebias} multiplies it by $h^2/2$, the bias grows without bound as the
fold is approached, and because $b = 0$ lies in the support of the estimate
(Lemma~\ref{lem:diverge}) the bias is unbounded at points the estimator is
expected to resolve.
\end{proof}

Lemma~\ref{lem:kdebias} explains an artefact we observed repeatedly and which
Section~\ref{sec:results} reports: the fold's apparent brightness and width in
a photon map track the kernel radius rather than the physics. Narrowing $h$
reduces the smoothing but raises the variance, since fewer photons fall inside
the kernel; the two effects pull against each other, and no $h$ makes the fold
correct. This is the reason the photon-map caustic in
Figure~\ref{fig:fourpanel} looks soft where ours is sharp, and it is not a
matter of insufficient photons.

Progressive photon mapping \cite{hachisuka2008}, and its stochastic variant
\cite{hachisuka2009}, shrink the radius as samples
accumulate, which makes the estimator consistent \emph{for the integral}. It
achieves this by letting $h \to 0$ while $N \to \infty$ at a controlled rate.
The bias term of \eqref{eq:kdebias} then vanishes in the limit, at the price of
the sample count and the accumulation time that Section~\ref{sec:current}
measures. The estimator becomes correct in the limit and remains wrong at any
finite budget, and for a fold the convergence is slow precisely because the
integrand is unbounded: the reported SPPM timings in this paper
(Table~\ref{tab:cost}) are the cost of pushing that limit far enough to see the
feature at all.

\subsection{Backward: connection, and the same factor}
\label{sec:backward}

Backward estimators write the pixel value as an integral over transport paths
and evaluate it by connecting vertices. A connection between a receiver vertex
and a light-carrying vertex carries the geometric term
\begin{equation}
G(\mathbf{x} \to \mathbf{u}) \;=\;
\frac{\big|\nrm_{\mathbf{u}} \cdot \omega_{\mathbf{u}}\big|}
{\|\mathbf{x} - \mathbf{u}\|^{2}} \cdot
\big|\detJ(\mathbf{x})\big|^{-1}
\label{eq:geometricterm}
\end{equation}
up to the conventions of the particular formulation. The last factor is the
same $1/|\detJ|$: the connection is undefined where the map is degenerate.
Backward methods therefore inherit the fold problem exactly, and the
literature's principal responses are to avoid connecting near degeneracies, to
solve for the connection on a manifold instead of sampling it
\cite{hanika2015,zeltner2020}, or to accept the noise and denoise it away.

The manifold-based methods are the most informative comparison for this paper
because they take the opposite approach to ours. Where we keep the geometric
quadrature everywhere and \emph{replace} the integrand inside the fold band,
manifold sampling keeps the geometric integrand and \emph{solves} for where the
degenerate connection must be. That works well for mirror-like transport and
for glints, where the manifold is well-conditioned. It does not address the
fold's divergence, because the divergence is a property of the integrand the
method still evaluates; near the fold the solution point is sought where the
Jacobian vanishes, which is the same place the weight blows up.

\subsection{The boundary of validity, stated once}
\label{sec:validity}

Collecting the two families, the situation is the following. Away from $\fold$
the map is a local diffeomorphism, $\detJ$ is bounded away from zero, and both
families are exact in the limit and cheap in practice; this is the regime in
which caustic rendering is a solved problem. On $\fold$ the density is
unbounded (Lemma~\ref{lem:diverge}), the forward estimator's bias is unbounded
for every bandwidth (Lemma~\ref{lem:kdebias}), and the backward estimator's
connection weight is unbounded. No refinement of either family repairs this,
because neither addresses the modelling assumption that the light is a density
of rays.

\begin{table}[tbp]
\centering
\caption{Where the standard families stand, by the analysis of this section.
``Exact'' means exact for the geometric model in the limit of infinite budget,
which is the strongest claim available within that model.}
\label{tab:families}
\small
% Column widths sum to 0.80\linewidth, not 0.86: with four columns the inter-column
% \tabcolsep padding adds 36 pt, and 0.86*240 + 36 = 242 pt against a 240 pt
% measure, which pushed the last column past the margin. The first column is
% wider than the rest because a family name is a single unbreakable word:
% "Manifold/vertex connection" has no break opportunity and overflowed a 46 pt
% column, which no font size fixes.
\begin{tabular}{@{}p{0.24\linewidth}p{0.17\linewidth}p{0.18\linewidth}p{0.21\linewidth}@{}}
\toprule
Family & Away from $\fold$ & On $\fold$ & Repair \\
\midrule
Photon map \cite{jensen1996} &
exact in limit, cheap &
bias $\sim h^2 b^{-5/2}$ unbounded &
none; narrower $h$ trades bias for variance \\
\addlinespace
SPPM \cite{hachisuka2008} &
exact in limit, expensive &
consistent but slow, since $h\to0$ is needed &
time alone \\
\addlinespace
Path tracing + NEE &
exact in limit, cheap for direct &
weight $\propto |\detJ|^{-1}$ unbounded &
variance only \\
\addlinespace
Manifold sampling \cite{hanika2015,zeltner2020} &
excellent, well conditioned &
solves at the degeneracy &
none; targets a different failure \\
\addlinespace
Wave-optical BRDF \cite{yan2018,cuypers2012} &
changes the model & changes the model &
not applicable to a transmissive fold \\
\bottomrule
\end{tabular}
\end{table}

\subsection{Where this sits in the transport literature}
\label{sec:related}

It is worth separating three things this paper borrows from, because a reader
may reasonably ask why a wave-optical correction is being attached to estimators
from a purely geometric literature.

The first is the \emph{estimator}. The push-forward of \eqref{eq:irradiance} is
what photon mapping estimates by density \cite{jensen1996}, what progressive
variants estimate with a shrinking radius \cite{hachisuka2008,hachisuka2009},
and what bidirectional and vertex-connection methods estimate by connecting
vertices rather than by storing samples \cite{veach1997}. Path-guiding and
manifold methods attack the same integral by changing \emph{where} the samples
go \cite{muller2017,hanika2015,zeltner2020}. Our correction changes none of
this. It changes the integrand inside a band and is therefore orthogonal to
every one of those improvements: a better sampler in the band still estimates an
unbounded function.

The second is the \emph{Jacobian}. Propagating irradiance through a scene by its
determinant is classical \cite{arvo1994}, and the rank deficiency that concerns
us is the same rank deficiency that appears in those formulations as a
degenerate geometric term. Treating the degeneracy as a modelling boundary
rather than as a numerical nuisance is the position this paper takes.

The third is the \emph{diffraction}. Uniform diffraction theories patch a
geometric solution with a correction near a shadow boundary
\cite{kouyoumjian1974}; catastrophe optics gives the corresponding local forms
for caustics \cite{berry1980,airy1838}; and graphics has applied wave theory to
reflectance \cite{cuypers2012,yan2018} and to propagation through media
\cite{musbach2013}. The band-local correction here is that same patch, restricted
to a fold and applied inside a renderer's estimator rather than to a field
everywhere.

Placing the method against all three at once also makes its scope clear. It
inherits the estimator's sampling behaviour away from the band, it inherits the
Jacobian's degeneracy as the thing it must detect, and it inherits the fold's
normal form as the only catastrophe it can currently correct. The parts of the
literature it does not touch are the parts that solve different problems.

\subsection{Wave optics already in graphics}
\label{sec:priorwave}

Treating light as a wave is not new to graphics, and the existing work divides
into two groups that are both different from what we do.

The first group modifies the reflectance model so that a surface, or a
microsurface, reflects according to wave theory. Cuypers et al.\
\cite{cuypers2012} give a diffraction-aware BRDF for a surface with a periodic
structure, and Yan et al.\ \cite{yan2018} render specular microgeometry with
wave optics, treating a rough surface as a collection of diffracting patches.
Both replace the reflection model and keep the transport geometric. Neither
addresses refraction through a transparent body, which is where a water-surface
caustic lives.

The second group models propagation through a medium with a wave equation
rather than with rays. Musbach et al.\ \cite{musbach2013} solve for the field
through a medium and produce diffraction effects in a rendering context. This is
closer in spirit to our treatment of the fold, and it is considerably more
expensive, because it propagates a field through the whole domain rather than
only near the degeneracy. The classical continuous-mathematics counterpart
is the uniform theory of diffraction \cite{kouyoumjian1974}, which patches a
geometric solution with a diffraction correction near a shadow boundary; the
band correction of this paper is the same manoeuvre applied to a caustic rather
than to an edge.

Our position sits between the two. We keep the geometric transport and the
geometric quadrature, which is where the efficiency is, and pay for wave optics
only inside a detected band, which is where the geometry is wrong. The cost of
that decision is a detector and a conserved-flux bookkeeping step; the benefit
is that the expensive part is confined to a set of small measure
(Section~\ref{sec:ours-cost}).
